\begin{sidewaystable}
\centering
\caption{Les hypothèses diagnostiques}
\label{tab:hypothesesDiagnostiques}
\begin{tabular} {|>{\raggedright\arraybackslash}p{0.1\textwidth}|>{\raggedright\arraybackslash}p{0.2\textwidth}|>{\raggedright\arraybackslash}p{0.2\textwidth}|>{\raggedright\arraybackslash}p{0.2\textwidth}|>{\raggedright\arraybackslash}p{0.2\textwidth}|}
        \toprule
        Diagnostic & Définition & Symptômes & Signes cliniques & Remarques \\
        \midrule
        \hline
        Kyste hémorragique du corps jaune & Kyste fonctionnel, dépendant du cycle menstruel. & Dysménorrhée avec douleur unilatérale lors de la phase lutéale. & \vspace{-3mm} \begin{itemize}[label=\textbullet]
            \item Palpation abdominale : sensible ;
            \item Toucher vaginale : masse annexielle palpée, mobilisation utérine douloureuse.
        \end{itemize}
        & Rupture du kyste \rightarrow \hspace{5pt} défense abdominale, signe du cri du Douglas. \\
        \hline
        Endométriome & Endomètre ectopique au niveau ovarien. & Dysménorrhée, douleurs pelviennes chroniques, dyspareunie profonde, troubles de la fertilité. & \vspace{-3mm} \begin{itemize}[label=\textbullet]
            \item Palpation abdominale : sensible ;
            \item Toucher vaginale : masse annexielle palpable, mobilisation utérine douloureuse (adhérences suspectées), cul-de-sac du Douglas sensible.
        \end{itemize}
        & Les symptômes associés à d’autres systèmes sont observés (nycturie, dyschésie, trouble du transit, \ldots). Cependant, aucune explication scientifique ne parvient à établir un lien de causalité entre les symptômes et le kyste. \\
        \hline
        Kyste dermoïde & Tératome, tumeur bénigne d'origine embryonnaire, fréquente chez les femmes entre $16$ et $55$ ans. & Asymptomatique ou légère douleur pelvienne, métrorragie, dyspareunie et symptômes liés à une compression éventuelle. & \vspace{-3mm} \begin{itemize}[label=\textbullet]
            \item Palpation abdominale : normal ;
            \item Toucher vaginale : masse annexielle peut-être palpable. 
        \end{itemize}
        & Torsion ovarienne = urgence gynécologique \rightarrow \hspace{5pt} douleur abdominale aiguë et intense.  \\
        \hline
        Kyste para-ovarien & Kyste bénin avec pour origine le mésosaplinx. & Asymptomatique ou symptômes liés à une compression éventuelle. & \vspace{-3mm} \begin{itemize}[label=\textbullet]
            \item Palpation abdominale : normal ;
            \item Toucher vaginale : masse annexielle peut-être palpable. 
        \end{itemize}
        &  \\
        \hline
        \vspace{-3mm} Tumeur ovarienne & \vspace{-3mm} Cancer ovarien malin. Facteurs favorisants : G$0$P$0$, ménarche précoce, cycle menstruel spontané, \ldots & \vspace{-3mm} \parbox{0.2\textwidth}{Stade précoce : paucisymptomatique.\\
        Stade avancé : AEG, perte de poids, fatigue, \ldots} 
        & \vspace{-3mm} \begin{itemize}[label=\textbullet]
            \item Palpation abdominale : signes de carcinose (ventre de batracien, \ldots);
            \item toucher vaginale : masse annexielle palpable voire d'autres masses selon le stade d'extension.
        \end{itemize} & \vspace{-3mm} Ce diagnostic doit être exclu vu son pronostic dramatique. \\
        \hline
        \end{tabular}
\end{sidewaystable}